% ===========================================================================
\section{Bounded harmonic functions}\label{sec:harmonic}
% ===========================================================================
Source: \cite{TAMSa}. A function $u$ defined in $\D$ has the \emph{asymptotic value} $a$ at $\zeta\in C$ if $u\to a$ along some arc
at $\zeta$, and the \emph{general limit} $a$ at $\zeta$ if $u(z)\to a$ as $z\to\zeta$, $z\in\D$. $P_r(\theta)=(1-r^2)/(1-2r\cos\theta+r^2)$
is the Poisson kernel, $\Delta(S,T)=\inf\{\abs{s-t+2m\pi}:s\in S,t\in T,m\in\Z\}$ for $S,T\subseteq\R$, and $\expi S=\{e^{it}:t\in S\}$.

\begin{fact}[Poisson integrals {\cite[pp.~32--38]{Hoffman}}]\label{fact:poisson}
\lean{TJK.Harmonic.poisson_bounded_harmonic, TJK.Harmonic.poisson_tendsto_of_continuousAt,
  TJK.Harmonic.hp_radial_limit_poisson}
\source{Hoffman, Banach spaces of analytic functions, pp.~32--38; Mathlib has the Poisson kernel (Analysis/Complex/Poisson)}
(i) If $\varphi$ is a bounded measurable real function on $C$, its Poisson integral
$P[\varphi](re^{i\theta})=\frac1{2\pi}\int_{-\pi}^{\pi}\varphi(e^{it})P_r(\theta-t)\,dt$ is a bounded harmonic function in $\D$ with
$\abs{P[\varphi]}\le\sup\abs\varphi$.
(ii) If $\varphi\in L^1(C)$ is defined on a set $S$ of full measure, $\zeta_0\in S$, and $\varphi|_S$ is continuous at $\zeta_0$, then
$P[\varphi]$ has general limit $\varphi(\zeta_0)$ at $\zeta_0$.
(iii) If $u$ is harmonic in $\D$ and $\{\int_{-\pi}^\pi\abs{u(re^{i\theta})}^pd\theta:0\le r<1\}$ is bounded for some $p>1$, then the
radial limit of $u$ exists at almost every point of $C$, the function $\psi$ so defined almost everywhere is in $L^1(C)$, and
$u=P[\psi]$.
\end{fact}

\begin{fact}[Jordan curve theorem and accessibility]\label{fact:jordan-access}
\lean{TJK.Harmonic.jordan_accessible}
\source{Jordan curve theorem; Newman, Elements, corollary on p.~164}
If $\Gamma\subseteq\C$ is a Jordan curve, $\C\setminus\Gamma$ has exactly two components, a bounded one $B$ and an unbounded one,
each with frontier $\Gamma$, and every point $p\in\Gamma$ is accessible from $B$: there is a simple arc $\gamma$ with endpoint $p$ and
$\gamma\setminus\{p\}\subseteq B$.
\end{fact}

\begin{lemma}[{\cite[Lemma]{TAMSa}}]\label{lem:ha-lemma}
\lean{TJK.Harmonic.accessible_avoiding_domains}
\uses{fact:jordan-access}
Let $(I_n)$ be pairwise disjoint nondegenerate closed arcs of $C$ and $(\Omega_n)$ Jordan domains in $\D$ with
$\overline{\Omega_n}\cap C=I_n$, $\overline{\Omega_n}\cap\overline{\Omega_m}=\emptyset$ for $n\ne m$, and $\diam\Omega_n\to0$. Then every point of
$C\setminus\bigcup_nI_n^*$ is accessible by a simple arc in $\D\setminus\bigcup_n\overline{\Omega_n}$.
\end{lemma}
\begin{proof}
$\partial\Omega_n$ is a Jordan curve containing $I_n$; let $\Gamma_n=\partial\Omega_n\setminus I_n^*$, an arc with the endpoints of $I_n$, and
$\phi_n:I_n\to\Gamma_n$ a homeomorphism fixing the endpoints. Let $\Gamma=(C\setminus\bigcup_nI_n^*)\cup\bigcup_n\Gamma_n$ and define $\phi:C\to\Gamma$
by $\phi=\phi_n$ on $I_n$ and $\phi=\mathrm{id}$ elsewhere; let $\phi_p$ use only $\phi_1,\dots,\phi_p$. Each $\phi_p$ is continuous, and
$\abs{\phi_p-\phi}\le\sup_{n>p}\diam\Omega_n\to0$, so $\phi$ is continuous; it is injective and $C$ is compact, so $\Gamma$ is a Jordan
curve. Each $\Omega_n$ misses $\Gamma$; a small disk about an interior point of $I_n$ misses $\Gamma$ and meets both $\Omega_n$ and
$\C\setminus\overline\D$, which lies in the unbounded component; so $\Omega_n$ lies in the unbounded component. The points of
$C\setminus\Gamma=\bigcup I_n^*$ also lie in the unbounded component (they are limits of points outside $\overline\D$ near which
$\Gamma$ is absent), and $\overline{\Omega_n}\subseteq\Omega_n\cup\Gamma_n\cup I_n$. Hence the bounded component $B$
(Fact~\ref{fact:jordan-access}) lies in $\D\setminus\bigcup_n\overline{\Omega_n}$, and each point of $C\setminus\bigcup I_n^*\subseteq\Gamma$
is accessible from $B$.
\end{proof}

\begin{lemma}[general limits along an arc]\label{lem:general-limit-endpoints}
\lean{TJK.Harmonic.asymptotic_at_endpoint}
\uses{def:ccset}
Let $u:\D\to\R$ have general limit $c$ at every point of an open arc $J$ of $C$. Then $u$ has the asymptotic value $c$ at each endpoint
of $J$.
\end{lemma}
\begin{proof}
Let $J=\expi(t_0,t_1)$ and $\tau_k=t_0+(t_1-t_0)2^{-k-1}$. By compactness of $\expi[\tau_k,\tau_0]$ there is $\rho_k\in(0,1/k)$ with
$\abs{u(re^{it})-c}<1/k$ for $t\in[\tau_k,\tau_0]$ and $1-\rho_k<r<1$; we may take $\rho_k$ decreasing. Let $r(t)$ be continuous on
$(t_0,\tau_0]$, affine on each $[\tau_{k+1},\tau_k]$, with $r(\tau_k)=1-\rho_{k+1}/2$. Then $\{r(t)e^{it}\}\cup\{e^{it_0}\}$ is an arc at
$e^{it_0}$ along which $u\to c$.
\end{proof}

\begin{lemma}[the intervals]\label{lem:ha-sequences}
\lean{TJK.Harmonic.exists_intervals}
Let $(p_n)$ be a countable dense subset of $(-\pi,\pi)$. There are natural numbers $n(1)<n(2)<\cdots$ and reals $\eps_k$ such that, with
$\alpha(k)=p_{n(k)}-\pi\eps_k^2/2$ and $\beta(k)=p_{n(k)}+\pi\eps_k^2/2$:
I.~$0<\eps_k\le2^{-k-1}$; II.~$-\pi<\alpha(k)<\beta(k)<\pi$; III.~the intervals $[\alpha(k),\beta(k)]$ are pairwise disjoint;
IV.~$\eps_k\le\bigl(\Delta([\alpha(h),\beta(h)],[\alpha(k),\beta(k)])/4\bigr)^4$ for $h<k$; V.~$p_h\in\bigcup_{i\le k}[\alpha(i),\beta(i)]$
for $h\le n(k)$.
\end{lemma}
\begin{proof}
Let $n(1)=1$ and choose $\eps_1$ satisfying I, II. Given the first $q$ terms, let $n(q+1)$ be the least $n$ with
$p_n\notin\bigcup_{i\le q}[\alpha(i),\beta(i)]$ (it exists because this finite union of closed intervals of total length $<2\pi$ cannot
contain a dense set; and $n(q+1)>n(q)$ by V), and choose $\eps_{q+1}>0$ with I, II and
$\frac\pi2\eps_{q+1}^2+4\eps_{q+1}^{1/4}\le\frac12\min_{k\le q}\Delta(p_{n(q+1)},[\alpha(k),\beta(k)])$. Then
$\Delta([\alpha(q+1),\beta(q+1)],[\alpha(k),\beta(k)])\ge4\eps_{q+1}^{1/4}>0$, which gives III and IV, and V holds by the minimality
of $n(q+1)$.
\end{proof}

\begin{lemma}[estimates]\label{lem:ha-estimates}
\lean{TJK.Harmonic.omega_jordan, TJK.Harmonic.omega_closure_disjoint, TJK.Harmonic.diam_omega_tendsto,
  TJK.Harmonic.fn_le_eps, TJK.Harmonic.fn_le_eps_sqrt, TJK.Harmonic.fn_tendsto_zero}
\uses{lem:ha-sequences, fact:poisson}
With the data of Lemma~\ref{lem:ha-sequences}, let $I_n=\expi[\alpha(n),\beta(n)]$, $\chi_n$ its indicator, $f_n=P[\chi_n]$, and
$\Omega_n=\{re^{i\theta}:1-\eps_n<r<1,\ \Delta(\theta,[\alpha(n),\beta(n)])<(1-r)^{1/4}\}$. Then $\Omega_n$ is a Jordan domain with
$\overline{\Omega_n}\cap C=I_n$; $\overline{\Omega_n}\cap\overline{\Omega_m}=\emptyset$ for $n\ne m$; $\diam\Omega_n\to0$; and
\begin{enumerate}[label=(\arabic*)]
\item $0\le f_n\le\eps_n$ on $\D\setminus\Omega_n$;
\item $f_n(re^{i\theta})\le\eps_n\sqrt{1-r}$ when $r>1-\eps_n$ and $\Delta(\theta,[\alpha(n),\beta(n)])\ge(1-r)^{1/4}$;
\item $f_n\to0$ along every arc at a point of $C$ that does not meet $\Omega_n$.
\end{enumerate}
\end{lemma}
\begin{proof}
$\Omega_n$ is the image of $(0,1)\times(1-\eps_n,1)$ under $(s,r)\mapsto r\exp\bigl(i[\alpha(n)-(1-r)^{1/4}+s(\beta(n)-\alpha(n)+2(1-r)^{1/4})]\bigr)$,
which extends to a homeomorphism of the closed square onto $\overline{\Omega_n}$ taking the edge $r=1$ onto $I_n$. A common point of
$\overline{\Omega_n}$ and $\overline{\Omega_m}$, $n<m$, would give $\Delta([\alpha(n),\beta(n)],[\alpha(m),\beta(m)])\le2\eps_m^{1/4}$, against IV.
$\diam\Omega_n\le\eps_n+\pi\eps_n^2+2\eps_n^{1/4}$.

For $re^{i\theta}\notin\Omega_n$ either $r\le1-\eps_n$ or ($r>1-\eps_n$ and $\Delta(\theta,[\alpha(n),\beta(n)])\ge(1-r)^{1/4}$). In the first case
$P_r\le(1+r)/(1-r)\le2/\eps_n$ and $f_n\le\frac1{2\pi}\cdot\pi\eps_n^2\cdot\frac2{\eps_n}=\eps_n$. In the second case, with $v=(1-r)^{1/4}<1$, for
$t\in[\alpha(n),\beta(n)]$ we have $\cos(\theta-t)\le\cos v\le1-\frac{v^2}2+\frac{v^4}{24}\le1-\frac14\sqrt{1-r}$; as $r>\frac34$,
\[
f_n\le\frac1{2\pi}\cdot\pi\eps_n^2\cdot\frac{1-r^2}{\frac12r\sqrt{1-r}}\le\frac83\eps_n^2\sqrt{1-r}\le\eps_n\sqrt{1-r}\le\eps_n,
\]
since $\eps_n\le\frac14$. This proves (1) and (2); (3) follows from (2) because $r\to1$ along such an arc.
\end{proof}

\begin{theorem}[{\cite{TAMSa}}; answer to {\cite[Problem 5]{BP}}]\label{thm:ha-example}
\lean{TJK.Harmonic.exists_bounded_harmonic_boundaryFn_not_baireOne}
\uses{lem:ha-estimates, lem:ha-sequences, lem:ha-lemma, lem:general-limit-endpoints, fact:poisson,
  cor:baire-continuity-point, lem:baire-precomp, def:ccset}
There are a bounded harmonic function $f$ in $\D$ and a boundary function $\varphi:C\to\R$ for $f$ that is not of Baire class $\le1$.
\end{theorem}
\begin{proof}
With the notation of Lemma~\ref{lem:ha-estimates}, let $\chi=\sum_n\chi_n$ (the indicator of $\bigcup_nI_n$) and $f=P[\chi]=\sum_nf_n$
(monotone convergence); $f$ is bounded harmonic (Fact~\ref{fact:poisson}(i)).

$\chi$ is continuous with value $1$ on each open arc $I_n^*$, so $f$ has general limit $1$ there (Fact~\ref{fact:poisson}(ii)) and, by
Lemma~\ref{lem:general-limit-endpoints}, asymptotic value $1$ at $e^{i\alpha(n)}$ and $e^{i\beta(n)}$.

Let $\zeta\in C\setminus\bigcup_nI_n^*$. By Lemma~\ref{lem:ha-lemma} there is an arc $\gamma$ at $\zeta$ meeting no $\Omega_n$. For $m\in\N$ and
$z\in\gamma$, by (1) and I, $\abs{f(z)}\le\sum_{n\le m}f_n(z)+\sum_{n>m}\eps_n\le\sum_{n\le m}f_n(z)+2^{-m-1}$, and by (3)
$\limsup\abs{f}\le2^{-m-1}$ along $\gamma$. So $f$ has asymptotic value $0$ at $\zeta$; in particular the endpoints of the $I_n$ are
ambiguous points. Let $A=\{e^{i\alpha(n)}\}$, $B=\{e^{i\beta(n)}\}$, and define $\varphi=1$ on $\bigcup_nI_n^*\cup B$ and $\varphi=0$ on the rest of
$C$. Then $\varphi$ is a boundary function for $f$.

By V, every $p_h$ lies in $\bigcup_nI_n$, so $\bigcup_nI_n$ is dense in $C$. Given $e^{i\alpha(n)}$ and a short open arc $O$ ending at it on
the side away from $I_n$, some $I_m$ meets $O$; since $I_m\cap I_n=\emptyset$, $e^{i\beta(m)}\in O$; and some $I_{m'}$ meets the open arc between
$e^{i\beta(m)}$ and $e^{i\alpha(n)}$, so lies inside it and $e^{i\alpha(m')}\in O$. Arguing likewise at points of $B$, each point of $A\cup B$
is a limit of other points of $A$ and of points of $B$. Hence $P=\cl(A\cup B)$ is a nonempty closed (so complete) set in which both $A$
and $B$ are dense. As $\varphi=0$ on $A$ and $\varphi=1$ on $B$, $\varphi|_P$ is continuous at no point. If $\varphi$ were of Baire class $\le1$, so
would be $\varphi|_P$, contradicting Corollary~\ref{cor:baire-continuity-point}.
\end{proof}

\begin{theorem}[{\cite[Theorem]{TAMSa}}]\label{thm:ha-theorem}
\lean{TJK.Harmonic.isMeagre_asymptotic_diff_generalLimit}
\uses{fact:poisson, cor:bf-honorary, thm:kuratowski, lem:meager-subspace}
Let $u$ be harmonic in $\D$ with $\{\int\abs{u(re^{i\theta})}^pd\theta:0\le r<1\}$ bounded for some $p>1$. Let $A$ be the set of points of $C$
at which $u$ has an asymptotic value and $G$ the set at which $u$ has a general limit. Then $A\setminus G$ is of the first category.
\end{theorem}
\begin{proof}
Let $S$ be the set where the radial limit $\psi$ exists; $S\subseteq A$, $C\setminus S$ is null, and $u=P[\psi]$ (Fact~\ref{fact:poisson}(iii)).
Extend $\psi|_S$ to a boundary function $\varphi_1$ on $A$. By Corollary~\ref{cor:bf-honorary} there is $\varphi\in\mathcal B_1(A,\R)$ with
$\varphi=\varphi_1$ off a countable set; then $\varphi=\psi$ almost everywhere and $u=P[\varphi]$. By Theorem~\ref{thm:kuratowski} there is $K\subseteq A$ with
$A\setminus K$ of the first category in $A$, hence in $C$ (Lemma~\ref{lem:meager-subspace}), such that $\varphi|_A$ is continuous at each point of
$K$. By Fact~\ref{fact:poisson}(ii), $K\subseteq G$; so $A\setminus G\subseteq A\setminus K$.
\end{proof}

\begin{corollary}[{\cite[Corollary]{TAMSa}}]\label{cor:ha-corollary}
\lean{TJK.Harmonic.isMeagre_discontinuity_boundaryFn}
\uses{thm:ha-theorem, def:ccset}
Under the hypothesis of Theorem~\ref{thm:ha-theorem}, if $\varphi$ is a boundary function for $u$ defined on $E\subseteq C$, then the set of
points of $E$ at which $\varphi$ is discontinuous is of the first category in $C$.
\end{corollary}
\begin{proof}
At a point $\zeta\in E\cap G$ with general limit $L$: if $\abs{u-L}<\eps$ on $\D\cap S(\delta,\zeta)$, then every $x'\in E\cap S(\delta,\zeta)$ has
$\abs{\varphi(x')-L}\le\eps$, since $\varphi(x')$ is a limit along an arc whose tail lies in $S(\delta,\zeta)$. So $\varphi$ is continuous on $E\cap G$, and
its discontinuities lie in $A\setminus G$.
\end{proof}

\begin{corollary}[analytic functions; {\cite{TAMSa}}]\label{cor:ha-analytic}
\lean{TJK.Harmonic.isMeagre_asymptotic_diff_generalLimit_holomorphic}
\uses{thm:ha-theorem, cor:ha-corollary}
Theorem~\ref{thm:ha-theorem} and Corollary~\ref{cor:ha-corollary} hold for holomorphic $F$ in $\D$ with bounded $p$-means, $p>1$.
\end{corollary}
\begin{proof}
Apply them to $\operatorname{Re}F$ and $\operatorname{Im}F$, whose $p$-means are bounded by those of $F$: if $F$ has an asymptotic value at $x$ so
do both parts, and if both have general limits so does $F$.
\end{proof}

\begin{remark}
The example of \cite[Theorem 6]{BP} shows that the hypothesis on the integral means cannot be omitted.
\end{remark}
